% captions.tex: hand-edited source of truth for the results-section figure titles
% and captions. One ptitle and one pcaption block per figure; the name matches the
% \pvtitle{name} / \pvcaption{name} references in section_source.tex (and the
% fig:<name> label). Titles are rendered by the TeX (as the bold lead-in of the
% caption), not baked into the PNGs. Blocks may use \pv{key} references: data
% values come from generated reporting fragments, including assessment-access
% paper_values.tex; display conventions come from figmeta.tex (figures.jl).
% Convention: captions describe elements and derivations only, no interpretation.

\begin{ptitle}{position}
Bridging position and bridging behavior
\end{ptitle}
\begin{pcaption}{position}
Points show ensemble late means; whiskers show pointwise \pv{centralIntervalPercent}\% Monte Carlo intervals. A--C use neural-network learning; parameters not varied in each panel remain at baseline. Principal degree includes the broker tie when present; maximum degree is calculated within each period before averaging. Lines connect matching compositions at fixed turnover in A and C and serve as visual guides across learning restrictions in D.
\end{pcaption}

\begin{ptitle}{assessment_access}
Brokerage with access, assessment, or both
\end{ptitle}
\begin{pcaption}{assessment_access}
Parameters other than turnover remain at baseline. A--C show late-mean differences: adding assessment is full service minus access-only; adding access is full service minus assessment-only. In B, dashed lines connect contrasts within each regime; principal degree includes the broker tie when present. Zero turnover is plotted separately in A and C. D shows unsmoothed ensemble means. Bars and ribbons show pointwise \pv{aaIntervalPercent}\% Monte Carlo intervals, paired by seed in A--C.
\end{pcaption}

\begin{ptitle}{information}
Sources of the broker's advantage
\end{ptitle}
\begin{pcaption}{information}
Ensemble late means from the Ridge experiments. Rank correlations use the same
randomly sampled candidates for principals and the broker. In C--D, parameters
other than $\rho$ remain at baseline; grey lines connect the same regime across
versions. Whiskers show pointwise \pv{infoIntervalPercent}\% Monte Carlo intervals.
C uses within-seed broker-minus-principal differences; B and D average seed-level
net-output shares. B's right-hand percentages compare ensemble-mean total net
output with the reference version.
\end{pcaption}

\begin{ptitle}{dynamics}
Broker use, access, and assessment accuracy
\end{ptitle}
\begin{pcaption}{dynamics}
A: ensemble means after a \pv{rollWin}-observation rolling mean within each seed, with pointwise 95\% Monte Carlo intervals. B--D: ensemble means for \pv{nRegimes} regimes; C--D use late means. Marginal curves show smoothed distributions of regime means, with equal weight per regime and common smoothing across learners.
\end{pcaption}
